% Zumdahl Chapter 3 test -- 20 MCQ + 2 FRQ. 50 min.
\documentclass{apchem-exam}

\apcunitword{Chapter}
\apcunit{3}
\apcunittitle{Stoichiometry}
\apcdoctitle{Chapter Test}
\apcsubtitle{Zumdahl \S3.1--3.11 \textbullet\ 50 minutes \textbullet\ periodic table provided}

\begin{document}
\apctitleblock

\examsection{I}{Multiple Choice}{20 questions \textbullet\ 30 minutes \textbullet\ 1 point each}{\nocalculator}

% ---- moles / molar mass ----------------------------------------------------
\begin{mcq}
  The molar mass of \ce{Ca(OH)2} is closest to
  \choices
    \choice{\SI{57}{\gram\per\mole}}
    \correct{\SI{74}{\gram\per\mole}}
    \choice{\SI{58}{\gram\per\mole}}
    \choice{\SI{90}{\gram\per\mole}}
\end{mcq}
\why{$40.08 + 2(16.00 + 1.008) = 74.10$.}
\distractor{A}{counted only one OH group}

\begin{mcq}
  How many moles are in \SI{88.0}{\gram} of \ce{CO2}?
  \choices
    \choice{\SI{1.00}{\mole}}
    \correct{\SI{2.00}{\mole}}
    \choice{\SI{4.00}{\mole}}
    \choice{\SI{0.500}{\mole}}
\end{mcq}
\why{$88.0/44.01 = 2.00$.}

\begin{mcq}
  Which sample contains the largest number of atoms?
  \choices
    \choice{\SI{1.0}{\mole} of \ce{O2}}
    \choice{\SI{1.0}{\mole} of \ce{H2O}}
    \correct{\SI{1.0}{\mole} of \ce{C2H6}}
    \choice{\SI{1.0}{\mole} of \ce{NaCl}}
\end{mcq}
\why{8 atoms per molecule beats 3, 2, and 2.}

\begin{mcq}
  A \SI{0.500}{\mole} sample of a compound has a mass of \SI{29.0}{\gram}.
  Its molar mass is
  \choices
    \choice{\SI{14.5}{\gram\per\mole}}
    \choice{\SI{29.0}{\gram\per\mole}}
    \correct{\SI{58.0}{\gram\per\mole}}
    \choice{\SI{116}{\gram\per\mole}}
\end{mcq}
\why{$M = m/n = 29.0/0.500$.}
\distractor{A}{multiplied rather than divided}

\begin{mcq}
  How many oxygen atoms are in \SI{0.500}{\mole} of \ce{Al2(SO4)3}?
  \choices
    \choice{$3.01\times10^{23}$}
    \choice{$1.81\times10^{24}$}
    \correct{$3.61\times10^{24}$}
    \choice{$7.23\times10^{24}$}
\end{mcq}
\why{12 O per formula unit: $0.500 \times 12 \times 6.022\times10^{23}$.}
\distractor{B}{used 4 oxygens instead of 12}

% ---- percent composition / formulas ----------------------------------------
\begin{mcq}
  The mass percent of nitrogen in \ce{NH4NO3} ($M = 80.05$) is closest to
  \choices
    \choice{17.5\%}
    \correct{35.0\%}
    \choice{28.0\%}
    \choice{63.0\%}
\end{mcq}
\why{Two N atoms: $28.02/80.05 = 0.350$.}
\distractor{A}{counted only one nitrogen}

\begin{mcq}
  A compound is 40.0\% C, 6.7\% H, and 53.3\% O by mass. Its empirical
  formula is
  \choices
    \choice{\ce{C2H4O2}}
    \correct{\ce{CH2O}}
    \choice{\ce{CHO}}
    \choice{\ce{C6H12O6}}
\end{mcq}
\why{Mole ratio $3.33 : 6.65 : 3.33 = 1{:}2{:}1$.}
\distractor{D}{that's a molecular formula, not the simplest ratio}

\begin{mcq}
  A compound has the empirical formula \ce{CH2} and a molar mass of
  \SI{70.1}{\gram\per\mole}. Its molecular formula is
  \choices
    \choice{\ce{C2H4}}
    \choice{\ce{C4H8}}
    \correct{\ce{C5H10}}
    \choice{\ce{C6H12}}
\end{mcq}
\why{$70.1/14.03 = 5$.}

\begin{mcq}
  In determining an empirical formula, a student obtains the mole ratio
  \ce{C}:\ce{O} $= 1 : 2.5$. The correct next step is to
  \choices
    \choice{round 2.5 to 3}
    \choice{round 2.5 to 2}
    \correct{multiply both values by 2}
    \choice{divide both values by 2.5}
\end{mcq}
\why{A .5 ratio means the true subscripts are double; rounding changes the
compound.}

\begin{mcq}
  In combustion analysis of a compound containing C, H, and O, the mass of
  oxygen in the sample is found by
  \choices
    \choice{measuring the \ce{O2} consumed}
    \choice{converting the \ce{CO2} mass to oxygen}
    \correct{subtracting the C and H masses from the sample mass}
    \choice{measuring the \ce{H2O} produced}
\end{mcq}
\why{Sample oxygen is indistinguishable from the \ce{O2} used to burn it, so
it must be obtained by difference.}

% ---- balancing -------------------------------------------------------------
\begin{mcq}
  When \ce{C4H10 + O2 -> CO2 + H2O} is balanced with the smallest whole-number
  coefficients, the coefficient of \ce{O2} is
  \choices
    \choice{6}
    \choice{10}
    \correct{13}
    \choice{8}
\end{mcq}
\why{\ce{2 C4H10 + 13 O2 -> 8 CO2 + 10 H2O}: $16+10 = 26$ O atoms.}

\begin{mcq}
  Balancing a chemical equation is justified by the law of
  \choices
    \correct{conservation of mass}
    \choice{multiple proportions}
    \choice{definite proportion}
    \choice{constant volume}
\end{mcq}
\why{Atoms are conserved, so mass is conserved.}

\begin{mcq}
  Which action is NEVER permitted when balancing an equation?
  \choices
    \choice{changing a coefficient}
    \correct{changing a subscript}
    \choice{multiplying all coefficients by 2}
    \choice{adding a coefficient of 1}
\end{mcq}
\why{Subscripts define the substance; changing them makes a different
compound.}

% ---- stoichiometry ---------------------------------------------------------
\begin{mcq}
  For \ce{N2 + 3 H2 -> 2 NH3}, how many moles of \ce{H2} are needed to react
  completely with \SI{4.0}{\mole} of \ce{N2}?
  \choices
    \choice{\SI{4.0}{\mole}}
    \choice{\SI{8.0}{\mole}}
    \correct{\SI{12}{\mole}}
    \choice{\SI{1.3}{\mole}}
\end{mcq}
\why{3:1 ratio.}
\distractor{D}{divided by 3 instead of multiplying}

\begin{mcq}
  For \ce{2 KClO3 -> 2 KCl + 3 O2}, \SI{0.400}{\mole} of \ce{KClO3} produces
  how many moles of \ce{O2}?
  \choices
    \choice{\SI{0.267}{\mole}}
    \choice{\SI{0.400}{\mole}}
    \correct{\SI{0.600}{\mole}}
    \choice{\SI{1.20}{\mole}}
\end{mcq}
\why{$0.400 \times 3/2$.}

\begin{mcq}
  In a mass-to-mass stoichiometry problem, the mole ratio from the balanced
  equation is used to convert
  \choices
    \choice{grams of A to moles of A}
    \correct{moles of A to moles of B}
    \choice{moles of B to grams of B}
    \choice{grams of A directly to grams of B}
\end{mcq}
\why{Coefficients relate moles only; molar masses do the gram conversions.}

% ---- limiting reactant / yield ---------------------------------------------
\begin{mcq}
  For \ce{2 H2 + O2 -> 2 H2O}, a mixture of \SI{4.0}{\mole} \ce{H2} and
  \SI{3.0}{\mole} \ce{O2} is ignited. The limiting reactant is
  \choices
    \correct{\ce{H2}}
    \choice{\ce{O2}}
    \choice{neither --- they are in exact ratio}
    \choice{cannot be determined}
\end{mcq}
\why{\ce{H2} needs only 2.0 mol \ce{O2}; oxygen is in excess.}
\distractor{B}{picks the smaller starting amount without using the ratio}

\begin{mcq}
  In the reaction above, how much \ce{O2} remains after the reaction is
  complete?
  \choices
    \choice{\SI{0}{\mole}}
    \correct{\SI{1.0}{\mole}}
    \choice{\SI{2.0}{\mole}}
    \choice{\SI{3.0}{\mole}}
\end{mcq}
\why{$3.0 - 2.0 = 1.0$ mol left.}

\begin{mcq}
  The theoretical yield of a reaction must always be calculated from
  \choices
    \choice{the reactant present in the greatest mass}
    \choice{the reactant with the largest molar mass}
    \correct{the limiting reactant}
    \choice{the average of all reactants}
\end{mcq}
\why{Product stops forming when the limiting reactant is consumed.}

\begin{mcq}
  A reaction with a theoretical yield of \SI{25.0}{\gram} gives an actual
  yield of \SI{20.0}{\gram}. The percent yield is
  \choices
    \choice{5.00\%}
    \choice{125\%}
    \correct{80.0\%}
    \choice{45.0\%}
\end{mcq}
\why{$20.0/25.0 \times 100$.}
\distractor{B}{inverted the ratio}

\examsection{II}{Free Response}{2 questions \textbullet\ 20 minutes}{\calculatorok}

% ---- FRQ 1 -----------------------------------------------------------------
\frq{short}{4}{Zumdahl \S3.6--3.7 \textbullet\ composition to formula}

A compound contains only carbon, hydrogen, and oxygen. Analysis shows it is
54.5\% C and 9.10\% H by mass, and its molar mass is
\SI{88.11}{\gram\per\mole}.

\begin{parts}
  \item Determine the mass percent of oxygen.
        \answerlines[1]{$100 - 54.5 - 9.10 = 36.4\%$}
  \item Determine the empirical formula. Show your mole work.
        \workspace[2.8cm]
        \keyonly{\begin{apcanswer}C: $54.5/12.01 = 4.54$;
        H: $9.10/1.008 = 9.03$; O: $36.4/16.00 = 2.28$. Divide by 2.28:
        $1.99 : 3.96 : 1.00 \to$ \ce{C2H4O}.\end{apcanswer}}
  \item Determine the molecular formula.
        \answerlines[2]{$M_{\text{emp}} = 44.05$; $88.11/44.05 = 2.00$;
        molecular formula \ce{C4H8O2}.}
\end{parts}

\begin{rubric}
  \rpoint{1}{(a) 36.4\% oxygen}
  \rpoint{1}{(b) converts all three percentages to moles correctly}
  \rpoint{1}{(b) \ce{C2H4O} after dividing by the smallest}
  \rpoint{1}{(c) \ce{C4H8O2} with the molar-mass ratio shown}
\end{rubric}

% ---- FRQ 2 -----------------------------------------------------------------
\frq{short}{4}{Zumdahl \S3.10--3.11 \textbullet\ limiting reactant and yield}

Aluminum reacts with chlorine: \ce{2 Al + 3 Cl2 -> 2 AlCl3}.
A mixture of \SI{27.0}{\gram} Al and \SI{142}{\gram} \ce{Cl2} is allowed to
react to completion.

\begin{parts}
  \item Determine the limiting reactant, showing the comparison that
        justifies your choice. \workspace[2.8cm]
        \keyonly{\begin{apcanswer}Al: $27.0/26.98 = \SI{1.00}{\mole} \to
        \SI{1.00}{\mole}$ \ce{AlCl3}. \ce{Cl2}: $142/70.90 =
        \SI{2.00}{\mole} \to \tfrac23(2.00) = \SI{1.33}{\mole}$ \ce{AlCl3}.
        Al gives less $\to$ Al is limiting.\end{apcanswer}}
  \item Calculate the theoretical yield of \ce{AlCl3} in grams.
        \answerlines[2]{$\SI{1.00}{\mole} \times \SI{133.33}{\gram\per\mole}
        = \SI{133}{\gram}$.}
  \item If \SI{112}{\gram} of \ce{AlCl3} is actually isolated, calculate the
        percent yield. \answerlines[1]{$112/133 \times 100 = 84.2\%$}
  \item State one reason a percent yield is typically less than 100\%.
        \answerlines[1]{Product lost during transfer or purification;
        incomplete reaction; competing side reactions.}
\end{parts}

\begin{rubric}
  \rpoint{1}{(a) both reactants converted to moles AND compared using the
    coefficient ratio --- comparing raw moles alone earns 0}
  \rpoint{1}{(a) identifies Al as limiting}
  \rpoint{1}{(b) \SI{133}{\gram} with work}
  \rpoint{1}{(c) 84.2\% AND (d) any valid physical reason}
\end{rubric}

\examtotal

\end{document}
